% LaTeX Canvas: Chapter Outline
\chapter{Scripting}
\label{ch:scripts-vs-notebooks}

\section*{Purpose}
As students progress, the limiting factor is rarely ``can you write code''; it is whether your work can be reused, rerun, and explained. Scripts complement notebooks by making work \emph{repeatable} from the command line and easier to automate, test, and share. This chapter teaches how to write simple scripts, import your own code into notebooks, translate notebooks into scripts, pass parameters from the command line, and make deliberate notebook-versus-script choices.

\section*{Learning objectives}
By the end of this chapter, a reader should be able to:
\begin{enumerate}
\item Write a Python script that can be run from the command line and imported as a module.
\item Organize reusable code into functions and (optionally) a small \texttt{src/} module.
\item Import custom scripts into a Jupyter notebook and reload changes during iteration.
\item Convert notebooks to scripts (and scripts to notebooks) using standard tools.
\item Pass parameters into scripts using a command-line interface (CLI) pattern.
\item Explain trade-offs: when a notebook is the right tool and when a script is.
\item Use a ``notebook as narrative, script as engine'' pattern to keep work reproducible.
\end{enumerate}

\section*{Running theme: separate \emph{logic} from \emph{presentation}}
The logic (data loading, cleaning, analysis) should live in importable functions. The notebook (or script) should orchestrate and explain.

\section{Mental models and vocabulary}
\subsection{Script, module, package}
\begin{description}
\item[Script] A file you run (e.g., \texttt{python analyze.py}).
\item[Module] A file you import (e.g., \texttt{import analysis}).
\item[Package] A directory of modules (typically with \texttt{\_\_init\_\_.py}) installed or importable.
\end{description}

\subsection{Entry point vs library code}
\begin{itemize}
\item \textbf{Entry point:} parsing parameters, calling functions, writing outputs.
\item \textbf{Library code:} reusable functions/classes that do the work.
\end{itemize}

\subsection{Working directory and paths}
\begin{itemize}
\item Scripts and notebooks depend on the working directory for relative paths.
\item Habit: define a project root and use consistent relative paths inside it.
\end{itemize}

\section{Writing scripts: the essentials}
\subsection{A minimal script structure}
\begin{itemize}
\item Imports
\item Constants/config
\item Functions
\item \texttt{main()} function
\item \texttt{if \_\_name\_\_ == "\_\_main\_\_":} block
\end{itemize}

\subsection{Why \texttt{main()} matters}
\begin{itemize}
\item Enables clean separation between code that runs on import vs code that runs on execution.
\item Makes the file usable both as a script and as an importable module.
\end{itemize}

\subsection{Inputs, outputs, and side effects}
\begin{itemize}
\item Inputs: file paths, URLs, parameters.
\item Outputs: files (CSV, figures), printed summaries, logs.
\item Prefer writing outputs to named folders (e.g., \texttt{data/processed}, \texttt{figures/}).
\end{itemize}

\subsection{Print vs logging (intro level)}
\begin{itemize}
\item Use print for small, human-facing status.
\item Introduce logging as a structured alternative for larger scripts.
\end{itemize}

\section{Organizing reusable code: from one file to \texttt{src/}}
\subsection{From monolith to functions}
\begin{itemize}
\item Identify repeated blocks and extract them into functions.
\item Make functions pure when possible: inputs in, outputs out.
\end{itemize}

\subsection{A simple project layout}
\begin{itemize}
\item \texttt{notebooks/} (narrative and exploration)
\item \texttt{src/} (reusable code)
\item \texttt{scripts/} (CLI entry points)
\item \texttt{data/raw}, \texttt{data/processed}, \texttt{figures/}
\item \texttt{README.md} (how to run)
\end{itemize}

\subsection{Import paths without pain (student-safe guidance)}
\begin{itemize}
\item Prefer running from the project root so imports and relative paths behave.
\item Avoid copying code between notebooks; import instead.
\item If imports are confusing, treat it as a project-structure problem, not a ``magic command'' problem.
\end{itemize}

\section{Loading custom scripts into notebooks}
\subsection{The basic pattern: import your code}
\begin{itemize}
\item Place reusable functions in \texttt{src/} (or a clearly named module file).
\item Import them in the notebook and call them like normal functions.
\item Keep notebooks thin: analysis orchestration + narrative.
\end{itemize}

\subsection{Iterating on imported code: reloading}
\begin{itemize}
\item Notebooks keep a long-lived kernel; imports do not automatically update.
\item Teach a safe workflow:
\begin{enumerate}
\item Edit code in the editor.
\item Reload the module (conceptual introduction to \texttt{importlib.reload}).
\item Re-run a small test cell.
\end{enumerate}
\end{itemize}

\subsection{Alternative: run a script from a notebook}
\begin{itemize}
\item Use notebook mechanisms to execute a script when appropriate.
\item Warn: running external processes can create confusing state splits (kernel state vs process state).
\end{itemize}

\subsection{Notebook ``smoke test'' cell}
\begin{itemize}
\item A small cell near the top that:
\begin{itemize}
\item prints working directory
\item checks required files exist
\item imports your \texttt{src/} modules
\item prints key versions
\end{itemize}
\end{itemize}

\section{Passing parameters from the command line}
\subsection{Why parameterize scripts}
\begin{itemize}
\item Make the same analysis run on different datasets or time ranges.
\item Enable automation (scheduled runs, batch processing).
\item Reduce copy/paste ``variant'' scripts.
\end{itemize}

\subsection{Three levels of parameterization}
\begin{description}
\item[Level 1: constants] edit values in the script (fast, not scalable).
\item[Level 2: config files] store parameters in a small text file (scalable for teams).
\item[Level 3: CLI arguments] pass parameters at run time (best for automation).
\end{description}

\subsection{The CLI pattern (student baseline)}
\begin{itemize}
\item Define required positional arguments (e.g., input file).
\item Define optional flags (e.g., \texttt{--out}, \texttt{--seed}, \texttt{--limit}).
\item Provide clear help text and defaults.
\end{itemize}

\subsection{Validating inputs}
\begin{itemize}
\item Check that input paths exist.
\item Validate parameter ranges (e.g., nonnegative integers).
\item Fail fast with helpful error messages.
\end{itemize}

\subsection{Reproducibility and auditability}
\begin{itemize}
\item Print or log the resolved parameters at start.
\item Write outputs to a deterministic location based on parameters.
\item Record the command used to run the script in \texttt{README.md} or an output metadata file.
\end{itemize}

\section{Translating notebooks into scripts (and vice versa)}
\subsection{Why convert}
\begin{itemize}
\item Scripts are easier to run in batch jobs and automation.
\item Scripts produce cleaner diffs in version control.
\item Converting can help reveal hidden-state problems.
\end{itemize}

\subsection{One-way conversion: notebook to script}
\begin{itemize}
\item Convert an \texttt{.ipynb} into a \texttt{.py} for review, reuse, or batch runs.
\item Treat the output as a starting point: refactor into functions.
\end{itemize}

\subsection{Round-trip conversion: paired notebooks}
\begin{itemize}
\item Pair an \texttt{.ipynb} with a text representation (e.g., percent script) so you can edit in a normal editor.
\item Benefit: better diffs, easier merges.
\end{itemize}

\subsection{Parameterizing notebooks (bridge to automation)}
\begin{itemize}
\item Some workflows keep notebooks but run them repeatedly with different parameters.
\item Introduce the concept of tagging a ``parameters'' cell and executing programmatically.
\end{itemize}

\section{Notebook versus script: trade-offs and decision rules}
\subsection{Notebooks are best for}
\begin{itemize}
\item Exploration and rapid iteration.
\item Teaching, explanation, and narrative analysis.
\item Sharing a computational story with plots and commentary.
\end{itemize}

\subsection{Scripts are best for}
\begin{itemize}
\item Repeatable runs with parameters.
\item Automation (cron, schedulers, CI).
\item Long-running jobs and robust logging.
\item Clean version control diffs and code review.
\end{itemize}

\subsection{Common hybrid pattern (recommended)}
\begin{itemize}
\item \textbf{Notebook:} narrative driver, figures, interpretation.
\item \textbf{src/:} reusable functions.
\item \textbf{scripts/:} CLI entry points for batch runs.
\end{itemize}

\subsection{Smells that suggest ``move to a script''}
\begin{itemize}
\item You manually rerun the same notebook with small changes.
\item The notebook takes a long time and you want it to run unattended.
\item You need command-line parameters or scheduled runs.
\item You need robust logs, retries, or consistent outputs.
\end{itemize}

\subsection{Smells that suggest ``stay in a notebook''}
\begin{itemize}
\item The primary goal is interpretation and communication.
\item The analysis is exploratory and changing quickly.
\item You are teaching or documenting reasoning.
\end{itemize}

\section{Best practices: making scripts and notebooks play well together}
\subsection{Make data and paths explicit}
\begin{itemize}
\item Avoid absolute paths.
\item Use a consistent project root.
\item Add a small path-check function used by both notebooks and scripts.
\end{itemize}

\subsection{Minimize hidden state}
\begin{itemize}
\item In notebooks: restart-and-run-all checks.
\item In scripts: pure functions and explicit parameters.
\end{itemize}

\subsection{Document how to run}
\begin{itemize}
\item A \texttt{README.md} with one ``copy/paste to run'' command.
\item Record expected inputs, outputs, and runtime.
\end{itemize}

\subsection{Version control discipline}
\begin{itemize}
\item Keep notebooks tidy: clear noisy outputs; avoid huge embedded data.
\item Prefer paired notebook formats or conversion to scripts for review.
\end{itemize}

\section{Worked examples (outline)}
\subsection{Example 1: Turn notebook code into importable functions}
\begin{itemize}
\item Identify a repeated transform.
\item Extract into \texttt{src/cleaning.py}.
\item Import and call from the notebook.
\end{itemize}

\subsection{Example 2: Write a CLI script around your functions}
\begin{itemize}
\item Create \texttt{scripts/run\_cleaning.py}.
\item Parse \texttt{--input} and \texttt{--output}.
\item Write outputs and print a short run summary.
\end{itemize}

\subsection{Example 3: Debug ``it imports in the terminal but not in the notebook''}
\begin{itemize}
\item Diagnose kernel/environment mismatch.
\item Diagnose project root/working directory mismatch.
\item Fix by selecting the correct kernel and running from the project root.
\end{itemize}

\subsection{Example 4: Convert a notebook to a script and refactor}
\begin{itemize}
\item Convert \texttt{analysis.ipynb} to \texttt{analysis.py}.
\item Clean the result: remove interactive-only cells, extract functions.
\item Add a minimal \texttt{main()} and CLI args.
\end{itemize}

\subsection{Example 5: Parameterize a notebook for repeated execution (optional)}
\begin{itemize}
\item Add a parameters cell.
\item Execute with different parameter sets and save outputs.
\end{itemize}

\section{Templates}
\subsection{Template A: Minimal script skeleton}
\begin{verbatim}
"""One-sentence purpose.

Inputs:
Outputs:
How to run:
"""

from pathlib import Path

def main(...):
...

if **name** == "**main**":
main(...)
\end{verbatim}

\subsection{Template B: Minimal CLI skeleton (conceptual)}
\begin{verbatim}
import argparse

def parse_args():
p = argparse.ArgumentParser(...)
p.add_argument("--input", required=True)
p.add_argument("--output", required=True)
return p.parse_args()

def main():
args = parse_args()
...

if **name** == "**main**":
main()
\end{verbatim}

\subsection{Template C: Notebook that calls scriptable functions}
\begin{verbatim}

# 1) Purpose + imports

# 2) Parameters (as variables)

# 3) Call functions from src/

# 4) Save outputs

# 5) Interpretation

# 6) Restart-and-run-all proof

\end{verbatim}

\section{Exercises}
\begin{enumerate}
\item Write a script that loads a CSV and prints a short summary (rows, columns, missingness).
\item Refactor the script so the summary logic is in a function and the script is only an entry point.
\item Import that function into a notebook and use it on two datasets.
\item Add a CLI with \texttt{--input} and \texttt{--output} flags.
\item Convert a notebook to a script and identify at least three hidden-state problems you had to fix.
\item Write a short paragraph explaining whether your project should be a notebook, a script, or a hybrid---and why.
\end{enumerate}

\section{One-page checklist}
\begin{itemize}
\item My script has a \texttt{main()} and does not run heavy work on import.
\item Reusable logic lives in importable functions/modules.
\item I can import my own code into notebooks and reload changes when iterating.
\item I can run the same analysis via a CLI with parameters.
\item I know how to convert notebooks to scripts and use conversion to improve reproducibility.
\item I choose notebooks for narrative/exploration and scripts for repeatable/automated runs.
\end{itemize}

\appendix
\section{Quick reference: common tools (conceptual)}
\begin{itemize}
\item Notebook conversion to scripts and reports.
\item Paired notebook formats for better diffs.
\item Parameterized execution of notebooks for batch runs.
\end{itemize}
